\chapter{Propuesta del Proyecto}

Este capítulo presenta la propuesta del \textit{Trabajo Fin de Grado}, delineando la estrategia metodológica adoptada y las fases de ejecución del proyecto. Aquí se define el marco operativo del trabajo, justificando las decisiones técnicas y metodológicas clave para abordar los desafíos identificados en el despliegue de \textit{Large Language Models} (\textit{LLM}) en entornos de \textit{HPC} bajo principios de \textit{SRE}.

\section{Fases del Proyecto}
\noindent\textbf{Fase de Hardware:}
\begin{itemize}
    \item Montaje físico del servidor, instalación del sistema operativo.
    \item Instalación de controladores y librerías necesarias: \textit{NVIDIA Driver}, \textit{CUDA Toolkit} y soporte para contenedores con \textit{GPU}.
\end{itemize}
\textbf{Fase de Despliegue con Podman (GPU):}
\begin{itemize}
    \item Creación de contenedores con \textit{Podman} para desplegar un modelo \textit{LLM} de gran tamaño (ej. \textit{DeepSeek}) y un modelo más pequeño.
    \item Configuración del entorno para aprovechar la aceleración por \textit{GPU}.
    \item Ejecución de pruebas de rendimiento (\textit{benchmarks}): inferencia, consumo de \textit{VRAM}, latencias, y rendimiento con carga concurrente.
\end{itemize}
\newpage
\noindent\textbf{Fase de Evaluación con OpenShift Local (SRE)}
\begin{itemize}
    \item Instalación y configuración de \textit{OpenShift Local}.
    \item Despliegue de servicios ligeros desde URLs de \textit{GitHub}, simulando pipelines \textit{DevOps}.
    \item Pruebas de capacidades \textit{SRE}.
    %\begin{itemize}
    %    \item Auto-scaling basado en carga simulada.
    %    \item Auto-healing tras fallos intencionados.
    %    \item Observabilidad y monitorización de recursos.
    %    \item Gestión de actualizaciones y reinicios controlados.
    %\end{itemize}
\end{itemize}

\section{Herramientas de Gestión y Planificación Temporal}
Para asegurar una ejecución eficiente y organizada del proyecto, se emplearon varias herramientas de gestión, documentación y planificación temporal. La organización general del \textit{Trabajo Fin de Grado} se basó en las recomendaciones del libro \textit{Cómo Escribir Tu TFG o TFM de Ingeniería Informática y No Morir en el Intento} ~\cite{ref34}. Se utilizaron \textit{Notion} para la gestión de documentación y notas, \textit{Clockify}~\cite{ref35} para el seguimiento del tiempo, y un diagrama de \textit{Gantt} para la planificación y visualización de las tareas. Además, todo el código y los \textit{scripts} desarrollados se versionaron y alojaron en un repositorio de \textit{GitHub} \cite{ref17} , garantizando la trazabilidad y accesibilidad del trabajo técnico.

El periodo de trabajo abarcó del 7 de octubre de 2024 al 6 de abril de 2025~\cite{ref31} (registro compartido), registrándose aproximadamente 300 horas para el TFG, esta cifra se aproximó al recomendado en el libro mencionado~\cite{ref34}. Se debe añadir que las horas registradas son solo indicativas y no incluyeron las sesiones de tutoría semanales. Asimismo, la realización de prácticas extracurriculares (del 1 de febrero al 31 de mayo, 35 horas semanales) podría considerarse como tiempo trabajado al no ser convalidadas.

Se incluyó un diagrama de \textit{Gantt}, una herramienta fundamental para la planificación temporal. \textit{Notion} demostró ser una herramienta muy potente para la organización y la planificación de ideas, ofreciendo, por ejemplo, la visualización de diagramas de \textit{Gantt}. Este diagrama, luego generado con \textit{Plotly} a partir de la información recopilada en \textit{Notion}, visualizó la organización y secuencia de tareas, desde el desarrollo hasta la redacción de la memoria, y reflejó la metodología de trabajo. Su presencia fue importante para cumplir los requisitos de la memoria y facilitar la evaluación.

\noindent El trabajo abordó los siguientes tipos de tareas:
\begin{itemize}
    \item T1: La preparación e implementación de un servidor \textit{GPU}.
    \item T2: El despliegue y la configuración de \textit{OpenShift} (software de virtualización).
    \item T3: El desarrollo y la experimentación con Modelos de Lenguaje (\textit{ML}) para una demostración.
    \item T4: La integración de principios de \textit{Site Reliability Engineering} (\textit{SRE}).
     \item T5: Redacción de la memoria del TFG (incluyendo introducción y estado del arte).
\end{itemize}

\newpage
    \begin{table}[H]
        \centering
        \begin{tabular}{|p{6cm}|c|c|c|}
            \hline
            \textbf{Nombre} & \textbf{Tarea} & \textbf{Fecha de comienzo} & \textbf{Tiempo} \\
            \hline
            Revisión y corrección final de todo el TFG pasarlo a latex & T5 & 30/05/2024 & 17:11:47 \\
            Hacer benchmarks y su Redacción & T3 & 29/05/2024 & 9:16:14 \\
            Terminar Demo openshift y primer borrador & T3 & 27/05/2024 & 3:09:15 \\
            Reorganizar + terminar benchmark + openshift & T5 & 27/05/2024 & 8:00:17 \\
            Redacción de la sección de Openshift Cap 6 y redacción Benchmark & T5 & 24/05/2024 & 8:00:41 \\
            Empezar los benchmarks para LLMs & T3 & 23/05/2024 & 10:05:00 \\
            Demo para openshift II & T2 & 19/04/2024 & 8:05:23 \\
            Demo para openshift & T2 & 18/04/2024 & 7:57:56 \\
            Desplegar modelo pequeño y uno grande cap 5 II & T3 & 18/04/2024 & 8:58:00 \\
            Desplegar modelo pequeño y uno grande cap 5 & T3 & 18/04/2024 & 10:08:00 \\
            Corregir la redacción de todo lo llevo & T5 & 11/04/2024 & 4:00:40 \\
            Redacción Estado del Arte III & T5 & 04/04/2024 & 7:17:32 \\
            Redacción de Montar el servidor cap 4 & T5 & 30/03/2024 & 8:09:26 \\
            Redacción Propuesta & T5 & 29/02/2024 & 7:07:45 \\
            Corrección introducción & T5 & 03/03/2024 & 3:28:30 \\
            Redacción Estado del Arte II & T5 & 15/02/2024 & 9:21:00 \\
            Redacción Estado del Arte & T5 & 08/02/2024 & 10:10:47 \\
            Inicio Curso Openshift IV & T2 & 05/02/2024 & 8:04:53 \\
            Inicio Curso Openshift III & T2 & 25/01/2024 & 7:08:10 \\
            Curso Openshift I & T2 & 18/01/2024 & 9:04:00 \\
            Curso Openshift II & T2 & 12/01/2024 & 8:30:22 \\
            Ver el Servidor en persona & T1 & 11/01/2024 & 11:02:14 \\
            Redacción Introducción & T5 & 19/12/2024 & 22:03:00 \\
            Notas Building Secure \& Reliable Systems O'Reilly & T4 & 10/12/2024 & 19:10:00 \\
            Redefinir Objetivos & T5 & 08/12/2024 & 20:26:46 \\
            Lecturas libro SRE & T4 & 01/12/2024 & 1:10:04 \\
            Estudiar Software de Virtualización: OpenShift & T2 & 25/11/2024 & 4:04:36 \\
            Definición inicial de objetivos & T5 & 10/10/2024 & 22:23:36 \\
            Estudio libro ``cómo hacer un TFG sin morir en el intento'' & T5 & 15/10/2024 & 20:33:50 \\
            Configuración entornos (Github, Notion) & T2 & 10/10/2024 & 7:00:00 \\
            Decisión tema TFG & T5 & 08/10/2024 & 1:06:09 \\
            \hline
        \end{tabular}
        \caption{Tabla para hacer el diagrama de Gantt}
        \label{tab:registro-tareas}
    \end{table}
    \addcontentsline{lof}{table}{Registro de tareas y planificación temporal}


\begin{figure}[H]
    \centering
    \includegraphics[width=1\textwidth]{imagenes/image69.png}
    \caption{Diagrama de Gantt generado con Plotly}
    \label{fig:diagrama-gantt-plotly}
\end{figure}

\begin{figure}[H]
    \centering
    \includegraphics[width=1\textwidth]{imagenes/image70.png}
    \caption{Resumen de clockify}
    \label{fig:resumen-clockify}
\end{figure}

\newpage
\section{Presupuesto y Gestión de Costes}
El desarrollo del \textit{TFG} se ha ejecutado con una gestión de costes optimizada, lo que se traduce en una inversión económica mínima. Fue posible gracias al acceso gratuito a recursos proporcionados por la universidad.

Un punto importante en la reducción de costes fue la formación en \textit{OpenShift}, facilitada sin coste por \textit{Red Hat} a los estudiantes, como se dijo anteriormente. Esta formación no solo dio las bases para el uso de \textit{OpenShift}, sino que también eliminó la necesidad de licencias de software de pago gracias a que proveyó de un entorno virtual privado con \textit{OpenShift} para usar.

Además, la universidad puso a disposición del proyecto un servidor \textit{Dell PowerEdge R7615} equipado con una \textit{GPU L40S}. Este acceso a hardware de alto rendimiento, indispensable para la viabilidad del proyecto, eliminó así la necesidad de incurrir en gastos de adquisición o alquiler de equipos especializados.

Sin embargo, al analizar el proyecto desde la perspectiva de una empresa que deseara replicar una iniciativa similar, resulta imprescindible considerar el coste real del hardware y del trabajo técnico. El servidor \textit{Dell PowerEdge R7615}, según el catálogo oficial de Dell España, tiene un precio base de 7.686,56€ \cite{ref45}. La \textit{GPU NVIDIA L40S 48GB GDDR6}, esencial para tareas de inferencia y entrenamiento con modelos de \textit{IA}, ronda entre 7.000€ y 8.000€ \cite{ref46}, según distribuidores españoles como Ibertrónica. El coste conjunto estimado para el hardware oscila, por tanto, entre 14.686,56€ y 15.686,56€.

En cuanto al coste del trabajo técnico, el desarrollo implicó una dedicación de unas 300 horas, como se indicó previamente con \textit{Clockify} \cite{ref35}. Tomando como referencia las tarifas de mercado para un ingeniero informático titulado en España, que oscilan entre 15€ y 18€ por hora según fuentes como \textit{indeed.com} o \textit{Talent.com} \cite{ref47,ref48} (son portales de empleo que también ofrecen datos salariales, en base a sus usuarios) se ha calculado un coste estimado de 5.586€ (300 horas * 17,13€/hora). Esta cifra es coherente con los valores del sector.

A pesar de que el proyecto se ejecutó con un coste directo prácticamente nulo gracias al apoyo de la universidad, su valor real estimado en el mercado supera los 20.000€. Esta cifra resalta tanto la dificultad de este trabajo como la eficiencia en el uso de los recursos disponibles, ofreciendo una base sólida para evaluar la viabilidad económica de iniciativas similares en entornos empresariales.
%\section{Contribución del Proyecto}
%Este proyecto se enfoca en evaluar la viabilidad y eficiencia del despliegue de \textit{LLM} en servidores \textit{HPC} mediante contenedores compatibles con \textit{GPU}, además de demostrar capacidades clave de \textit{SRE} utilizando \textit{OpenShift Local}. Este enfoque dual permite tanto el trabajo con modelos reales en el hardware disponible como la familiarización con tecnologías de \textit{orquestación} y fiabilidad operativa en entornos simulados. El trabajo se ha desarrollado con una perspectiva crítica, priorizando la automatización, resiliencia y optimización del rendimiento para ofrecer una visión práctica sobre los desafíos y soluciones en el despliegue de \textit{IA}.

%El TFG también consolida la experiencia práctica en el manejo de servidores adquirida durante las prácticas extracurriculares realizadas en Alemania, complementando y profundizando los conocimientos obtenidos. Se documentará el proceso de aprendizaje, incluyendo los aciertos, desafíos y soluciones implementadas.

%Desde una perspectiva profesional, este trabajo permite afianzar habilidades en la configuración y optimización de infraestructuras \textit{HPC}, el despliegue de \textit{LLM} y la aplicación de principios \textit{SRE}, competencias altamente demandadas en el ámbito de la \textit{IA} y la computación en la nube. Personalmente, este proyecto ha representado un reto significativo que ha permitido integrar conocimientos teóricos y prácticos, fomentando la resolución de problemas complejos y la autonomía. En última instancia, la intención es ofrecer una referencia formativa útil para futuros estudiantes interesados en el campo de la computación de alto rendimiento y la fiabilidad operativa.

